% =====================================================================
\section{Память: RAM, Flash, PSRAM и NVS}
% =====================================================================
\lessonintro
  {выбирать место для данных и сохранять настройки так, чтобы не изнашивать Flash}
  {урок 7 (задачи и общий доступ), урок 9 (данные после сна)}
  {яркость и имя устройства переживают отключение питания}
  {ничего нового}

\subsection{У каждой памяти своя роль}

RAM хранит переменные и стек задач: она быстрая, но очищается при сбросе. Внешняя PSRAM даёт больше
места для больших буферов, кадров камеры и веб-страниц, но медленнее и есть не на каждой плате.
Flash содержит прошивку и переживает отключение питания. Внутри Flash выделен раздел NVS --- маленькое
журналируемое хранилище «ключ--значение» для настроек.

\begin{figure}[H]
\centering
\begin{tikzpicture}[x=1.0cm,font=\sffamily\scriptsize]
  \foreach \a/\b/\c/\t in {0/1.3/blkrtc/{загруз-\\чик},1.3/2.9/blkrtc/{таблица\\ разделов},2.9/4/blkmem/nvs,
                          4/5.2/blkmem/{ota-\\data},5.2/8.3/blkcpu/{app0 (ota\_0)\\ прошивка},8.3/11.4/blkcpu/{app1 (ota\_1)\\ прошивка},
                          11.4/14.3/blkper/{SPIFFS / FAT\\ файлы},14.3/15.4/gray!15/{core-\\dump}}
    \draw[fill=\c] (\a,0) rectangle (\b,1) node[pos=0.5,align=center]{\t};
  \draw[decorate,decoration={brace,mirror,amplitude=4pt}] (2.9,-0.1) -- (4,-0.1) node[midway,below=5pt,align=center]{настройки\\ (\code{Preferences})};
  \draw[decorate,decoration={brace,mirror,amplitude=4pt}] (5.2,-0.1) -- (11.4,-0.1) node[midway,below=5pt]{два слота для OTA-обновления (урок 13)};
  \node[anchor=south west] at (0,1.05) {0x0};
  \node[anchor=south east] at (15.4,1.05) {конец Flash};
\end{tikzpicture}
\caption{Разметка Flash по умолчанию в Arduino-ESP32 (схема, не в масштабе). Выбрать другую таблицу
разделов можно в \textsf{Tools → Partition Scheme}}
\end{figure}

\begin{center}\small
\begin{tabularx}{\textwidth}{@{}l X X@{}}
\toprule
Где & Подходит для & Не подходит для \\
\midrule
RAM & текущие измерения, очереди, стек & сохранения после питания \\
PSRAM & большой временный буфер & критичного к задержкам кода \\
RTC RAM & счётчик deep-sleep, короткое состояние & сохранения после обычного сброса \\
NVS во Flash & настройки, калибровка, пароль Wi-Fi & частой телеметрии и журнала \\
SD / сервер & большой журнал и файлы & мгновенного доступа без питания \\
\bottomrule
\end{tabularx}
\end{center}

\subsection{Настройка в NVS}

Класс \code{Preferences} в Arduino-ESP32 безопасно работает с NVS. Пространство имён \code{"app"}
отделяет ваши ключи от других частей прошивки; ключи делайте короткими и постоянными.

\codecap{Чтение значения по умолчанию и редкое сохранение}
\codeinput{l11_nvs}

\begin{caution}
Flash нельзя записывать на каждом проходе \code{loop()}. Ресурс ячеек ограничен: запись уровня
потенциометра сотни раз в секунду быстро превратит полезный раздел в источник ошибок. Сохраняйте
настройку по действию пользователя, после тайм-аута «ручка перестала двигаться» или с ограничением,
например, не чаще раза в минуту.
\end{caution}

\subsection{PSRAM и большие буферы}

Наличие PSRAM нужно проверить: \code{psramFound()} возвращает \code{true}, а размер ---
\code{ESP.getPsramSize()}. Большой буфер выделяют через \code{ps\_malloc()}, проверяют на \code{nullptr}
и освобождают \code{free()}. Не переносите туда бездумно всё подряд: стек задач и маленькие горячие
структуры лучше оставить в обычной RAM. Утечки в долгоживущем веб-сервере удобно ловить сравнением
\code{ESP.getFreeHeap()} до и после сотен запросов.

\subsection{Разбор ошибки: настройки «забываются»}

Сначала убедитесь, что \code{prefs.begin("app", false)} открыт в режиме записи, а ключ при чтении и
записи совпадает. Затем посмотрите на результат \code{put...}: он возвращает число записанных байтов.
Не вызывайте \code{clear()} при обычном старте --- это функция сброса настроек. Наконец, не путайте
NVS с RTC-памятью: \code{RTC\_DATA\_ATTR} переживает deep-sleep, но не отключение питания.

\begin{summary}
\begin{itemize}[leftmargin=*]
  \item RAM --- для работы сейчас, NVS --- для редких постоянных настроек, RTC RAM --- между снами.
  \item PSRAM полезна для больших буферов, но сначала проверяем её наличие и результат выделения.
  \item Редкая осмысленная запись бережёт Flash и делает поведение устройства предсказуемым.
\end{itemize}
\end{summary}

\begin{selfcheck}
\begin{enumerate}
  \item Где хранить счётчик пробуждений, пароль Wi-Fi и большой буфер изображения?
  \item Почему Flash нельзя записывать на каждом проходе \code{loop()}?
  \item Как проверить, что PSRAM есть и память действительно выделилась?
\end{enumerate}
\end{selfcheck}

\begin{tasks}
\begin{enumerate}
\item \lv{1} Сохраните в NVS число запусков и выведите его после перезагрузки.
\item \lv{1} Сбросьте плату питанием и убедитесь, что значение сохранилось. Чем NVS отличается от обычной переменной?
\item \lv{2} Сохраните последнее состояние светодиода (вкл/выкл и яркость) и восстановите его после включения.
\item \lv{2} Реализуйте команду \code{factory}: удалите пространство имён \code{app} после подтверждения.
\item \lv{3} Добавьте отложенную запись яркости через 3 секунды после последнего поворота ручки, чтобы не изнашивать Flash.
\item \lv{4} Выделите в PSRAM буфер на 1 МБ, заполните его и проверьте содержимое; сравните скорость с обычной кучей.
\end{enumerate}
\end{tasks}
